% 来源及取材范围见分题文件；此为整理者转录的正文，不是下载原件。
\paragraph{Setting (source \S4.4).}
\noindent\textit{以下背景与记号据所列原文章节摘编；不是逐字引文，也不是新增证明。}\par An $(n,d,\lambda)$-graph is a $d$-regular graph on $n$ vertices whose adjacency matrix has eigenvalues $d=\lambda_1\geq\cdots\geq\lambda_n$ and $\max\{|\lambda_2|,|\lambda_n|\}\leq\lambda$.
\begin{theorem}[Result proved in Proof 2 following Theorem 4.15]
\[\max\{|\lambda_2|,|\lambda_n|\}\geq2\sqrt{d-1}-o(1).\]
\end{theorem}
\begin{proof}[Proof 2 (slightly weaker result)]
We'll show that $\max\{|\lambda_2|,|\lambda_n|\}\geq2\sqrt{d-1}-o(1)$. This is an illustration of the trace method, also called the moment method. We have
\[\sum_{i=1}^n\lambda_i^{2k}=\operatorname{tr}(A^{2k}).\]
The right hand side is the number of closed walks of length $2k$ on $G$. Now, the number of closed walks of length $2k$ starting at a fixed vertex $v$ in a $d$-regular graph is at least the number of closed walks of length $2k$ starting at a fixed $v$ in an infinite $d$-regular tree. To see why this is true, note that given any walk on the infinite $d$-regular tree, we can walk in the same way on $G$ by assigning an orientation to each vertex. But $G$ may have more walks if it has cycles.

There are at least $C_k(d-1)^k$ closed walks of length $2k$ starting at a fixed $v$ in an infinite $d$-regular tree, where $C_k=\frac1{k+1}\binom{2k}{k}$ is the $k$-th Catalan number. Thus, the number of walks of length $2k$ on $G$ is at least $\frac{n}{k+1}\binom{2k}{k}(d-1)^k$. On the other hand,
\[d^{2k}+(n-1)\lambda^{2k}\geq\sum_{i=1}^n\lambda_i^{2k}.\]
Thus,
\[\lambda^{2k}\geq\frac1{k+1}\binom{2k}{k}(d-1)^k-\frac{d^{2k}}{n}.\]
The term $\frac1{k+1}\binom{2k}{k}$ is $(2-o(1))^{2k}$ as $k\to\infty$. Letting $k\to\infty$ and $k=o(\log n)$ as $n\to\infty$ gives us $\lambda\geq2\sqrt{d-1}-o(1)$.
\end{proof}
